\documentclass[11pt,a4paper]{article}

\usepackage[utf8]{inputenc}
\usepackage[margin=1in]{geometry}
\usepackage{amsmath,amssymb,amsfonts}
\usepackage{booktabs}
\usepackage{multirow}
\usepackage{graphicx}
\usepackage{hyperref}
\usepackage{cite}
\usepackage{microtype}
\usepackage{authblk}
\usepackage{tabularx}
\usepackage{enumitem}

\hypersetup{
    colorlinks=true,
    linkcolor=blue,
    citecolor=blue,
    urlcolor=blue
}
\providecommand{\doi}[1]{doi: \href{https://doi.org/#1}{#1}}

% ---------------------------------------------------------------
% Title
% ---------------------------------------------------------------
\title{\textbf{Race-Progression Finish Time Prediction in IRONMAN 70.3
Triathlons: Evaluating Discipline Information Value and Location
Transferability via Supervised Learning}}

% ---------------------------------------------------------------
% Authors
% ---------------------------------------------------------------
\author[1]{Truong Minh Khue}
\author[1]{Nguyen Song Thu}
\author[1]{Pham Phuong Thao}
\author[1]{Do Thi Gia Hoi}
\affil[1]{Department of Artificial Intelligence, FPT University
Da Nang, Vietnam}
\affil[ ]{\textit{fudaresearch@gmail.com}}

\date{}

\begin{document}

\maketitle

% ---------------------------------------------------------------
% Abstract
% ---------------------------------------------------------------
\begin{abstract}
Predicting athlete finish time in long-distance multi-sport events such
as the IRONMAN~70.3 triathlon (1.9\,km swim, 90\,km bike, 21.1\,km
run) supports in-race pacing strategy and event operations.
Although previous studies have investigated machine learning-based
performance prediction using athlete demographics, geographic
characteristics, and race split information, the predictive value of
progressively available race information and the robustness of
predictions across temporal periods and event locations warrant further
investigation.
We propose a sequential, race-progression supervised learning framework
evaluated at three observation checkpoints: $C_0$ (pre-race
registration data), $C_1$ (post-swim, adding the official swim split),
and $C_2$ (post-bike, adding the official bike split).
On a quality-controlled cohort of 1{,}152{,}545 individual finishes
(2004--2023; audited from 1{,}152{,}658 pre-QC eligible entries), we
benchmark Linear Regression (LR), Random Forest (RF), and an un-tuned
CatBoost reference model (120 iterations, seed 42) against three
domain-specific historical-median baselines.
Using out-of-period temporal evaluation---Phase~I: trained on
2004--2016, tested on 2019--2020; Phase~II: trained on 2004--2018,
tested on 2021--2023---we quantify the MAE reduction associated with
each sequential checkpoint.
\texttt{CatBoost\_ref} MAE decreases from 36.37\,min at $C_0$ to
28.43\,min at $C_1$ (a 7.94-min reduction with swim time) and further by
13.64\,min to 14.79\,min at $C_2$ in Phase~I; in Phase~II, MAE decreases
from 41.72\,min at $C_0$ to 33.45\,min at $C_1$ (an 8.27-min reduction)
and further by 17.67\,min to 15.79\,min at $C_2$, outperforming the
checkpoint-aware median-remaining-time baseline~($B_2$) by
6.10--7.36\,min at $C_1$ and $C_2$.
Further analyses are planned to examine geographic generalization via
location-based hierarchical modeling and to quantify prediction
uncertainty using Conformalized Quantile Regression.
\end{abstract}

\noindent\textbf{Keywords:} Sports analytics; finish-time prediction;
supervised learning; IRONMAN~70.3; temporal validation.

% ---------------------------------------------------------------
% 1. Introduction
% ---------------------------------------------------------------
\section{Introduction}

The IRONMAN~70.3---also called the Half-Ironman---is a widely
contested long-distance triathlon format comprising three consecutive
disciplines: a 1.9\,km (1.2-mile) open-water swim, a 90\,km (56-mile)
cycling stage, and a 21.1\,km (13.1-mile) half-marathon run, totaling
113.0\,km (70.3\,miles). Competition duration for age-group athletes
typically ranges from four to eight hours, making accurate finish-time
estimation valuable for strategic pacing, hydration planning, and
event operations.

From a statistical-learning perspective, finish-time prediction is
inherently sequential. Before the starting gun, analysts have access
only to static demographic and logistical attributes: athlete sex, age
bracket, country of origin, race venue, and calendar year. As the race
unfolds, each timing mat yields new course-split measurements.
Clearing the swim mat ($C_1$) reflects both an athlete's swimming
capability and race-day water conditions. Transitioning off the
bicycle ($C_2$) reflects cumulative elapsed time and pacing across
more than two-thirds of the total race distance, providing a strong
predictor of final performance.

\subsection{Literature Review and Related Work}

Three bodies of work directly motivate this study.

\begin{enumerate}[leftmargin=*]
  \item \textbf{Demographic and pre-race forecasting.}
  Thuany et al.~\cite{thuany2025modeling} analysed 823{,}464 race
  records from IRONMAN~70.3 competitions using tree-based regression
  algorithms (including Random Forest and CatBoost) trained on sex,
  age group, country, and event location. Their study showed that
  demographic and venue covariates capture baseline performance patterns
  at the pre-race stage, but it did not examine sequential in-race
  updating.

  \item \textbf{Discipline-level correlations.}
  Nikolaidis et al.~\cite{nikolaidis2023predicting} examined the
  statistical relationships between individual split times and overall
  finish time across 823{,}459 age-group IRONMAN~70.3 finishes using
  multiple regression, correlation, and ANOVA. They investigated how
  swimming, cycling, and running splits relate to overall finish time
  without establishing a single dominant discipline across age groups;
  however, their analysis focused on retrospective split associations
  rather than prospective out-of-period checkpoint forecasting.

  \item \textbf{Uncertainty quantification.}
  Romano, Patterson, and Cand\`{e}s~\cite{romano2019conformalized}
  introduced Conformalized Quantile Regression (CQR), which provides
  distribution-free prediction intervals with finite-sample coverage
  guarantees under exchangeability. In athletic performance modeling,
  point estimates alone are insufficient; a calibrated interval
  communicates the practical range within which an athlete's finish
  time may fall.
\end{enumerate}

\subsection{Initial Two-Checkpoint Experiment and Research Gaps}
\label{sec:tce}

To motivate the present design, we first conducted an exploratory
two-checkpoint experiment (hereafter TCE) using the integrated
multi-source dataset~\cite{kaggle_dataset} that underpins this study.
Table~\ref{tab:tce} summarises the configuration and reported outputs
of the TCE; the original notebook and summary metrics are archived in
the project repository~\cite{repo_artifact}.

\begin{table}[htbp]
\centering
\small
\caption{Two-checkpoint initial experiment (TCE) configuration and
results. Features for $C_0$: sex, age band, continent, athlete type,
calendar year. $C_2$ added swim time and bike time simultaneously
(no intermediate post-swim checkpoint). Data split: trained on
2004--2020 ($N = 814{,}715$); evaluated on 2021--2023
($N = 323{,}967$). Finish times were trimmed to full-sample P1/P99
prior to training and categorical encoders were fit before the
train--test split.}
\label{tab:tce}
\renewcommand{\arraystretch}{1.2}
\begin{tabularx}{\textwidth}{llcccc}
\toprule
\textbf{Checkpoint} & \textbf{Model} & $N_{\text{fit}}$ &
$N_{\text{test}}$ & \textbf{MAE (min)} & \textbf{RMSE (min)} \\
\midrule
$C_0$ (pre-race)  & Linear Regression & 814,715 & 323,967 & 41.44 & 51.30 \\
$C_0$ (pre-race)  & Random Forest     & 814,715 & 323,967 & 40.89 & 51.16 \\
$C_2$ (post-bike) & Linear Regression & 814,715 & 323,967 & 16.81 & 34.89 \\
$C_2$ (post-bike) & Random Forest     & 814,715 & 323,967 & 15.39 & 20.34 \\
\bottomrule
\end{tabularx}
\end{table}

Moving from $C_0$ to $C_2$ in the TCE, RF~MAE fell by 25.50\,min
(40.89 to 15.39\,min), a large apparent reduction. However, because
of preprocessing leakage (full-sample P1/P99 target trimming and
pre-split categorical encoding), the TCE serves strictly as a
motivational pilot rather than an unbiased benchmark. It also leaves
three methodological questions open:

\begin{itemize}[leftmargin=*]
  \item \textbf{Conflation of split-discipline value.} Because swim
  time and bike time were injected simultaneously, the 25.50-min
  reduction blends the contributions of both disciplines. It is unknown
  how much error reduction is achieved immediately after the swim
  split alone.

  \item \textbf{Absence of a checkpoint-aware baseline.}
  The TCE lacked a domain-specific split-aggregation sanity check
  (Baseline~2, defined in Section~\ref{sec:baselines}) that adds
  observed split times to demographic historical medians of remaining
  time. Without $B_2$, algorithmic learning cannot be separated from
  simple split arithmetic.

  \item \textbf{Unverified out-of-period and spatial generalization.}
  The TCE used a single cut (2004--2020 / 2021--2023) and omitted
  out-of-period validation within a homogeneous source as well as
  systematic venue transfer evaluation.
\end{itemize}

In response to these gaps, this study establishes a leakage-free,
three-checkpoint pipeline evaluated across two out-of-period temporal
phases on an identical paired cohort, benchmarks three predictive
models against three domain baselines, and defines planned extensions
for venue-level hierarchical modeling and prediction uncertainty.

% ---------------------------------------------------------------
% 2. Research Objectives and Questions
% ---------------------------------------------------------------
\section{Research Objectives and Questions}

We structure our overarching research program around three Research
Questions (RQs). In the current empirical evaluation, we fully answer
\textbf{RQ1}, while \textbf{RQ2} and \textbf{RQ3} are formulated as
methodological extensions for subsequent implementation.

\begin{itemize}[leftmargin=*]
  \item \textbf{RQ1 (Completed) -- Discipline incremental value and
  algorithmic benchmarking.}
  \begin{itemize}
    \item \textbf{RQ1a.} Upon completing the swim split ($C_1$) and
    subsequently the bike split ($C_2$), by how many minutes do MAE,
    RMSE, Bias, and $P_{90}$ change relative to the immediately
    preceding checkpoint, and how do learned models compare with the
    checkpoint-aware median-remaining-time baseline~$B_2$?
    \item \textbf{RQ1b.} At each identical checkpoint with identical
    feature access ($F_3$), how does an un-tuned CatBoost reference
    model (\texttt{CatBoost\_ref}) perform relative to Linear
    Regression (LR) and Random Forest (RF) across both homogeneous
    (Phase~I) and cross-source (Phase~II) evaluations?
  \end{itemize}

  \item \textbf{RQ2 (Planned) -- Spatial generalization and venue
  transferability.} Across future calendar years, how does
  incorporating host venue identity ($F_3$ versus $F_2$) influence
  error when evaluated separately on historically seen race locations
  versus completely novel venues, and can Location-based Hierarchical
  Regression improve generalization?

  \item \textbf{RQ3 (Planned) -- Uncertainty quantification.} At each
  checkpoint, what empirical coverage and interval width do historical
  residual intervals, CatBoost Quantile Regression, and CQR achieve
  relative to nominal 80\% and 90\% targets on reserved calibration
  and evaluation years?
\end{itemize}

% ---------------------------------------------------------------
% 3. Methodology
% ---------------------------------------------------------------
\section{Methodology}

Table~\ref{tab:notation} summarizes the checkpoint and feature-set
notation used throughout this paper.

\begin{table}[htbp]
\centering
\small
\caption{Summary of checkpoint and geographic feature notation.}
\label{tab:notation}
\renewcommand{\arraystretch}{1.15}
\begin{tabularx}{\textwidth}{l X}
\toprule
\textbf{Symbol} & \textbf{Definition} \\
\midrule
$C_0$ & Pre-race checkpoint: registration and event attributes only ($S_{0,i} = 0$) \\
$C_1$ & Post-swim checkpoint: $C_0$ attributes plus official swim split $W_i$ ($S_{1,i} = W_i$) \\
$C_2$ & Post-bike checkpoint: $C_1$ attributes plus official bike split $B_i$ ($S_{2,i} = W_i + B_i$) \\
$F_0$ & Core non-geographic features: \texttt{gender}, \texttt{ageband}, \texttt{eventyear} \\
$F_1$ & $F_0$ plus athlete continent (planned ablation) \\
$F_2$ & $F_0$ plus athlete country ISO code (planned ablation) \\
$F_3$ & $F_2$ plus host venue name \texttt{eventlocation} (primary feature set evaluated here) \\
\bottomrule
\end{tabularx}
\end{table}

\subsection{Study Cohort, Data Integration, and Quality Control}

The dataset integrates historical race records from Kaggle
(2004--2020) and endurance-data.com (2021--2023) into three linked
relational tables (\texttt{Athletes}, \texttt{Races}, and
\texttt{Results}) verified against an internal relational database
snapshot (\texttt{1,153,194} raw rows; exact zero-discrepancy match on all
1{,}152{,}658 eligible records across all identifiers, demographics, and
split times). To guarantee valid comparisons, the following ingestion
and quality-control (QC) rules were enforced:

\begin{enumerate}[leftmargin=*]
  \item \textbf{Population eligibility.} Individual amateur age-group
  triathletes (\texttt{AthleteType = Age-Group}) with sex recorded as
  M or~F, belonging to one of eleven standardized five-year brackets
  (18--24 through 70--74). \texttt{AgeBand} is extracted directly as
  the integer lower bound of each bracket (18, 25, \dots, 70).

  \item \textbf{Split validity and complete-case guarantee.} Official
  swim time~($W_i$), bike time~($B_i$), and total finish time~($Y_i$)
  must all be strictly positive and finite simultaneously at ingestion.
  Because all three timing fields are validated up front, no record is
  dropped between $C_0$, $C_1$, and $C_2$, guaranteeing an identical
  paired test sample across all three checkpoints.

  \item \textbf{String hygiene.} All CSVs were parsed with
  \texttt{keep\_default\_na=False} so Namibia's country code
  (\texttt{"NA"}) is preserved as a valid string rather than coerced to
  missing. Categorical fields (\texttt{gender}, \texttt{countryiso},
  \texttt{eventlocation}) had leading/trailing whitespace stripped,
  with empty strings mapped to an explicit \texttt{"\_\_MISSING\_\_"}
  token.

  \item \textbf{Uniform logical-consistency audit (2004--2023).} A
  logical consistency check verifying $Y_i > W_i + B_i$ was executed
  uniformly across all twenty calendar years (2004--2023). Zero
  violations occurred in 2004--2020 (all training partitions and
  Phase~I evaluation/calibration partitions). Across 2021--2023,
  113~records exhibited $Y_i \le W_i + B_i$ (58 in 2021, 46 in 2022,
  and 9 in 2023), indicating anomalous timing artifacts or unreported
  disqualifications. These 113 records were excluded from evaluation
  and final metric computation. No percentile trimming was applied to
  finish times.\footnote{Unlike the TCE, which trimmed finish times to
  full-sample P1/P99 and fit encoders before splitting, the present
  protocol retains the full untruncated distribution and fits all
  preprocessing strictly on training splits.}

  \item \textbf{Final analytical cohort.} Auditing 1{,}152{,}658
  eligible records and removing the 113 inconsistent 2021--2023 records
  retained \textbf{1{,}152{,}545} unique individual finishes spanning
  2004--2023 ($697{,}267 + 125{,}912 + 329{,}366$).
\end{enumerate}

\subsection{Race-Progression Checkpoints}

Finish-time prediction is formulated as a supervised regression task
whose target $Y_i \in \mathbb{R}^{+}$ is the official elapsed finish
time (stored in seconds, evaluated in minutes). Using the primary
feature representation $F_3$:

\begin{itemize}[leftmargin=*]
  \item \textbf{Checkpoint 0 ($C_0$, pre-race):}
  $X_{0,i} = [\text{Gender}_i,\; \text{AgeBand}_i,\;
  \text{CountryISO}_i,\; \text{EventLocation}_i,\;
  \text{EventYear}_i]$, with $S_{0,i} = 0$.

  \item \textbf{Checkpoint 1 ($C_1$, post-swim):}
  $X_{1,i} = [X_{0,i},\; W_i]$, with $S_{1,i} = W_i$.

  \item \textbf{Checkpoint 2 ($C_2$, post-bike):}
  $X_{2,i} = [X_{0,i},\; W_i,\; B_i]$, with $S_{2,i} = W_i + B_i$.
\end{itemize}

\noindent\textbf{Feature-leakage guardrail.} Running splits, final
rankings, transition times (T1 and T2), split ratios, and any
quantity derived from $Y_i$ are excluded from all feature sets.

\subsection{Temporal Partitioning}
\label{sec:temporal}

Models are evaluated in two strictly out-of-period temporal phases
with disjoint \texttt{raceid} sets between fit and test
(Table~\ref{tab:temporal_split}).

\begin{table}[htbp]
\centering
\small
\caption{Temporal partitioning protocol. Counts reflect post-QC
records; the 113 QC exclusions apply exclusively to the Phase~II test
partition (329{,}479 pre-QC $\rightarrow$ 329{,}366 post-QC).}
\label{tab:temporal_split}
\renewcommand{\arraystretch}{1.2}
\begin{tabularx}{\textwidth}{lcccc X}
\toprule
\textbf{Phase} & \textbf{Training} & \textbf{Reserved Calib.}
  & \textbf{Test} & $N_{\text{fit}}$ & $N_{\text{test}}$ (races) \\
\midrule
I -- Homogeneous & 2004--2016 & 2017--2018 & 2019--2020
  & 495{,}489 & 125{,}912 (127) \\
II -- Cross-source & 2004--2018 & 2019--2020 & 2021--2023
  & 697{,}267 & 329{,}366 (258) \\
\bottomrule
\end{tabularx}
\end{table}

\textbf{Phase~I (within-source temporal evaluation)} trains on Kaggle
records from 2004--2016 ($N = 495{,}489$) and tests on 2019--2020
($N = 125{,}912$; 127 races). Years 2017--2018 ($N = 201{,}778$) are
held out from training and reserved strictly for future interval
calibration.

\textbf{Phase~II (cross-period, cross-source evaluation)} trains on
2004--2018 ($N = 697{,}267$), reserves 2019--2020 for future
calibration, and evaluates on 2021--2023 records from Endurance-Data
($N = 329{,}366$ post-QC; 258 races). Because both time period and data
source change simultaneously, Phase~II measures combined temporal,
source-coverage, and event-composition shifts rather than pure
temporal drift alone.

\subsection{Domain Baselines}
\label{sec:baselines}

Three non-parametric baselines are fit exclusively on each training
partition:

\begin{itemize}[leftmargin=*]
  \item $B_0$ \textbf{(global median):}
  $\hat{Y}_{B_0,i} = \mathrm{Median}(Y_{\text{fit}})$.

  \item $B_1$ \textbf{(demographic-stratified median):}
  $\hat{Y}_{B_1,i} = \mathrm{Median}(Y_{\text{fit}} \mid
  \text{Gender}_i, \text{AgeBand}_i)$.

  \item $B_2$ \textbf{(checkpoint-aware median-remaining-time baseline):}
  Adds the known completed discipline split(s) $S_{t,i} \in \{0, W_i,
  W_i + B_i\}$ (strictly excluding T1 and T2, matching model inputs)
  to the training-set median of remaining time stratified solely by
  sex and age bracket (not conditioned on individual split magnitude):
  \begin{equation}
    \hat{Y}_{B_2,i} = S_{t,i}
    + \mathrm{Median}(Y_{\text{fit}} - S_t
    \mid \text{Gender}_i, \text{AgeBand}_i).
  \end{equation}
  If a ($\text{Gender}$, $\text{AgeBand}$) training cell has fewer
  than 30 observations, the estimate falls back to the sex-level
  median, and then to the global median. At $C_0$ ($S_{0,i} = 0$),
  $B_1$ and $B_2$ share the identical code path and outputs ($B_1 \equiv B_2$).
\end{itemize}

\subsection{Supervised Predictive Models and Implementation}

All experiments were executed in Python~3.12 (\texttt{numpy==2.3.5},
\texttt{pandas==2.2.3}, \texttt{scikit-learn==1.8.0},
\texttt{catboost==1.2.10}). Each model is fit independently for every
$(\text{phase}, \text{checkpoint})$ pair.

\subsubsection{Shared Preprocessing for LR and RF}
For both Linear Regression and Random Forest, categorical features
(\texttt{gender}, \texttt{countryiso}, \texttt{eventlocation}) are
encoded using \texttt{OneHotEncoder(handle\_unknown='ignore')}, which
learns category dictionaries strictly from the training split and maps
any unseen test category to an all-zero indicator vector (without
frequency-threshold pooling). Numeric features (\texttt{ageband},
\texttt{eventyear}, and active checkpoint splits) are standardized via
\texttt{StandardScaler} fit solely on the training split to ensure
numerical stability in sparse least squares.

\subsubsection{Linear Regression (LR)}
An additive main-effects Ordinary Least Squares model
(\texttt{LinearRegression(tol=1e-9)}) trained on the transformed design
matrix without explicit interaction terms.

\subsubsection{Random Forest (RF)}
An ensemble of 200 regression trees (\texttt{max\_depth = 12},
\texttt{min\_samples\_leaf = 50}, \texttt{random\_state = 42}). The RF
hyperparameters were retained from the TCE to provide a consistent tree
architecture; however, results are not directly comparable across the
two studies because the underlying datasets, preprocessing procedures,
feature sets, and evaluation protocols differ.

\subsubsection{CatBoost Reference Model (\texttt{CatBoost\_ref})}
CatBoost~\cite{prokhorenkova2018catboost} employs ordered boosting and
ordered categorical feature statistics designed to mitigate prediction
shift associated with target-dependent statistics. Native categorical
indices (\texttt{cat\_features = ['gender', 'countryiso',
'eventlocation']}) are passed directly alongside unscaled numeric
features. To provide a reproducible, un-tuned benchmark without
post-hoc selection, \texttt{CatBoost\_ref} uses a fixed configuration:
120 iterations, tree depth 6, learning rate $\eta = 0.08$, L2 leaf
regularization $\lambda = 3$, RMSE objective,
\texttt{one\_hot\_max\_size}~= 255, and random seed 42.

\subsubsection{Planned Extensions: Location-Based Hierarchical
Regression (HR) and CQR}
\label{sec:hr}
To address RQ2 and RQ3 in subsequent phases, we formulate: (i)~a
linear mixed-effects model $Y_{ij} = X_{ij}\boldsymbol{\beta} + u_j +
\epsilon_{ij}$ estimated via REML, where $u_j \sim \mathcal{N}(0,
\sigma_u^2)$ captures venue-specific adjustments with empirical Bayes
shrinkage and fallback $u_j = 0$ for unseen venues; and
(ii)~Conformalized Quantile Regression~\cite{romano2019conformalized}
calibrated on the reserved intervening years (2017--2018 for Phase~I;
2019--2020 for Phase~II) for 80\% and 90\% prediction intervals.
Neither HR nor CQR is evaluated in the current empirical tables.

\subsection{Evaluation Metrics}

Predicted finish times $\hat{Y}_i$ and ground-truth labels $Y_i$ are
compared in minutes across all test records:
\begin{enumerate}[leftmargin=*]
  \item \textbf{MAE} (primary metric):
  $\mathrm{MAE} = \frac{1}{N}\sum_{i=1}^{N}|\hat{Y}_i - Y_i|$.
  \item \textbf{RMSE}:
  $\mathrm{RMSE} = \sqrt{\frac{1}{N}\sum_{i=1}^{N}(\hat{Y}_i-Y_i)^2}$.
  \item \textbf{Bias}:
  $\mathrm{Bias} = \frac{1}{N}\sum_{i=1}^{N}(\hat{Y}_i - Y_i)$
  (positive values denote over-prediction of finish time).
  \item \textbf{$P_{90}$ absolute error:} the 90th percentile of
  $|\hat{Y}_i - Y_i|$.
  \item \textbf{Skill score} relative to $B_2$ at the same checkpoint:
  $\mathrm{Skill} = \bigl(1 - \frac{\mathrm{MAE}_{\mathrm{model}}}
  {\mathrm{MAE}_{B_2}}\bigr)\times 100\%$.
\end{enumerate}
Although race-cluster bootstrap resampling (10{,}000 resamples by
\texttt{raceid}) is planned for statistical testing of pairwise
$\Delta\mathrm{MAE}$ once full per-record prediction archives are
attached, the present study reports exact complete-cohort aggregate
metrics verified across all 32 evaluation runs.

% ---------------------------------------------------------------
% 4. Results
% ---------------------------------------------------------------
\section{Results}

Table~\ref{tab:all_results} presents the complete experimental metrics
across all 32 model runs spanning Phase~I ($N_{\text{test}} =
125{,}912$; 127 races) and Phase~II ($N_{\text{test}} = 329{,}366$; 258
races). Table~\ref{tab:transitions} summarizes checkpoint-to-checkpoint
MAE reductions and skill scores relative to $B_2$.

\begin{table}[htbp]
\centering
\scriptsize
\caption{Complete experimental results across Phase~I (train
2004--2016, test 2019--2020, $N_{\text{fit}}=495{,}489$,
$N_{\text{test}}=125{,}912$) and Phase~II (train 2004--2018, test
2021--2023, $N_{\text{fit}}=697{,}267$, $N_{\text{test}}=329{,}366$).
All errors are in minutes ($\text{Bias} = \hat{Y} - Y$). Bold marks the
lowest MAE within each $(\text{Phase}, \text{Checkpoint})$ cell.}
\label{tab:all_results}
\renewcommand{\arraystretch}{1.12}
\begin{tabular}{lllrrrr}
\toprule
\textbf{Phase} & \textbf{CP} & \textbf{Model} & \textbf{MAE (min)}
  & \textbf{RMSE (min)} & \textbf{Bias (min)} & \textbf{$P_{90}$ (min)} \\
\midrule
\multirow{16}{*}{I (2019--2020)}
  & $C_0$ & $B_0$                  & 40.03 & 50.38 & $-2.21$ & 80.75 \\
  & $C_0$ & $B_1$                  & 38.03 & 47.94 & $-1.60$ & 75.87 \\
  & $C_0$ & $B_2$                  & 38.03 & 47.94 & $-1.60$ & 75.87 \\
  & $C_0$ & LR                     & 36.73 & 45.78 & $+5.48$ & 74.52 \\
  & $C_0$ & \textbf{\texttt{CatBoost\_ref}} & \textbf{36.37} & \textbf{45.63} & $+1.32$ & \textbf{73.66} \\
  & $C_0$ & RF                     & 37.06 & 46.73 & $-0.36$ & 75.07 \\
\cmidrule{2-7}
  & $C_1$ & $B_0$                  & 40.03 & 50.38 & $-2.21$ & 80.75 \\
  & $C_1$ & $B_2$                  & 34.55 & 43.59 & $-1.64$ & 68.67 \\
  & $C_1$ & LR                     & 30.77 & 39.60 & $+7.25$ & 61.68 \\
  & $C_1$ & \textbf{\texttt{CatBoost\_ref}} & \textbf{28.43} & \textbf{35.96} & $+2.11$ & \textbf{56.88} \\
  & $C_1$ & RF                     & 28.66 & 36.48 & $+0.86$ & 57.52 \\
\cmidrule{2-7}
  & $C_2$ & $B_0$                  & 40.03 & 50.38 & $-2.21$ & 80.75 \\
  & $C_2$ & $B_2$                  & 21.23 & 27.20 & $-2.37$ & 41.77 \\
  & $C_2$ & LR                     & 14.89 & 19.31 & $+2.68$ & 31.22 \\
  & $C_2$ & \textbf{\texttt{CatBoost\_ref}} & \textbf{14.79} & \textbf{19.26} & $+0.91$ & \textbf{31.17} \\
  & $C_2$ & RF                     & 15.19 & 19.92 & $+0.13$ & 32.20 \\
\midrule
\multirow{16}{*}{II (2021--2023)}
  & $C_0$ & $B_0$                  & 45.35 & 57.15 & $-11.97$ & 94.15 \\
  & $C_0$ & $B_1$                  & 43.06 & 54.31 & $-10.26$ & 88.08 \\
  & $C_0$ & $B_2$                  & 43.06 & 54.31 & $-10.26$ & 88.08 \\
  & $C_0$ & \textbf{LR}            & \textbf{41.70} & \textbf{51.86} & $-0.47$ & \textbf{83.92} \\
  & $C_0$ & \texttt{CatBoost\_ref} & 41.72 & 52.24 & $-5.12$ & 84.61 \\
  & $C_0$ & RF                     & 42.06 & 52.70 & $-4.76$ & 85.29 \\
\cmidrule{2-7}
  & $C_1$ & $B_0$                  & 45.35 & 57.15 & $-11.97$ & 94.15 \\
  & $C_1$ & $B_2$                  & 39.55 & 50.14 & $-9.91$ & 81.18 \\
  & $C_1$ & LR                     & 34.43 & 44.12 & $-0.50$ & 68.91 \\
  & $C_1$ & \textbf{\texttt{CatBoost\_ref}} & \textbf{33.45} & \textbf{43.17} & $-4.43$ & \textbf{68.12} \\
  & $C_1$ & RF                     & 33.65 & 43.50 & $-4.70$ & 68.79 \\
\cmidrule{2-7}
  & $C_2$ & $B_0$                  & 45.35 & 57.15 & $-11.97$ & 94.15 \\
  & $C_2$ & $B_2$                  & 23.15 & 30.03 & $-7.77$ & 49.10 \\
  & $C_2$ & LR                     & 16.55 & 22.53 & $-0.81$ & 34.57 \\
  & $C_2$ & \textbf{\texttt{CatBoost\_ref}} & \textbf{15.79} & \textbf{21.02} & $-1.67$ & \textbf{33.33} \\
  & $C_2$ & RF                     & 15.81 & 21.04 & $-1.69$ & 33.57 \\
\bottomrule
\end{tabular}
\end{table}

\begin{table}[htbp]
\centering
\small
\caption{Checkpoint-to-checkpoint MAE reductions (min) and skill
scores relative to $B_2$ (computed from unrounded metrics).}
\label{tab:transitions}
\renewcommand{\arraystretch}{1.15}
\begin{tabular}{llcccccc}
\toprule
\textbf{Phase} & \textbf{Model}
  & $\Delta_{C_0\to C_1}$ & $\Delta_{C_1\to C_2}$ & $\Delta_{C_0\to C_2}$
  & \textbf{Skill ($C_0$)} & \textbf{Skill ($C_1$)} & \textbf{Skill ($C_2$)} \\
\midrule
\multirow{4}{*}{I}
  & $B_2$                  & $-3.48$ & $-13.32$ & $-16.80$ & ---      & ---       & ---       \\
  & LR                     & $-5.96$ & $-15.88$ & $-21.84$ & $+3.41\%$ & $+10.94\%$ & $+29.88\%$ \\
  & RF                     & $-8.40$ & $-13.47$ & $-21.87$ & $+2.54\%$ & $+17.04\%$ & $+28.45\%$ \\
  & \texttt{CatBoost\_ref} & $-7.94$ & $-13.64$ & $-21.58$ & $+4.36\%$ & $+17.70\%$ & $+30.33\%$ \\
\midrule
\multirow{4}{*}{II}
  & $B_2$                  & $-3.50$ & $-16.40$ & $-19.91$ & ---      & ---       & ---       \\
  & LR                     & $-7.27$ & $-17.88$ & $-25.15$ & $+3.16\%$ & $+12.96\%$ & $+28.53\%$ \\
  & RF                     & $-8.42$ & $-17.83$ & $-26.25$ & $+2.31\%$ & $+14.94\%$ & $+31.70\%$ \\
  & \texttt{CatBoost\_ref} & $-8.27$ & $-17.67$ & $-25.93$ & $+3.10\%$ & $+15.42\%$ & $+31.80\%$ \\
\bottomrule
\end{tabular}
\end{table}

\subsection{RQ1a: Incremental Value of Swim and Bike Splits vs.\ $B_2$}

Across both evaluation phases and every supervised model, forecast
errors decline monotonically from $C_0$ to $C_1$ and $C_2$
(Table~\ref{tab:all_results}):

\begin{itemize}[leftmargin=*]
  \item \textbf{Post-swim update ($C_0 \to C_1$).} Adding the official
  swim split alone yields a sizable MAE drop across all ML models:
  \texttt{CatBoost\_ref} improves by 7.94\,min in Phase~I
  (36.37 $\to$ 28.43\,min) and 8.27\,min in Phase~II
  (41.72 $\to$ 33.45\,min); RF drops by 8.40\,min (Phase~I) and
  8.42\,min (Phase~II). In contrast, the checkpoint-aware baseline
  $B_2$ improves by only 3.48\,min (Phase~I) and 3.50\,min (Phase~II).
  Consequently, \texttt{CatBoost\_ref} outperforms $B_2$ by 6.11\,min
  (Skill $+17.70\%$) in Phase~I and 6.10\,min (Skill $+15.42\%$) in
  Phase~II at $C_1$.

  \item \textbf{Post-bike update ($C_1 \to C_2$).} Incorporating the
  bike split produces an even larger absolute reduction:
  \texttt{CatBoost\_ref} MAE falls by an additional 13.64\,min to
  14.79\,min in Phase~I, and by 17.67\,min to 15.79\,min in Phase~II.
  The margin over $B_2$ widens to 6.44\,min (Skill $+30.33\%$) in
  Phase~I and 7.36\,min (Skill $+31.80\%$) in Phase~II.

  \item \textbf{Tail-error contraction ($P_{90}$).} From $C_0$ to
  $C_2$, 90th-percentile absolute error for \texttt{CatBoost\_ref}
  contracts by $57.69\%$ in Phase~I (73.66 $\to$ 31.17\,min) and by
  $60.61\%$ in Phase~II (84.61 $\to$ 33.33\,min).
\end{itemize}

\subsection{RQ1b: Algorithmic Comparison Across Checkpoints}

Comparing LR, RF, and \texttt{CatBoost\_ref} under identical $F_3$
information access reveals distinct regime behavior:

\begin{itemize}[leftmargin=*]
  \item \textbf{Pre-race ($C_0$).} Pre-race covariates explain modest
  variance compared to demographic medians, with ML models beating
  $B_1/B_2$ by only 0.97--1.66\,min in Phase~I and 1.00--1.36\,min in
  Phase~II. \texttt{CatBoost\_ref} achieves the lowest MAE in Phase~I
  (36.37\,min), whereas LR slightly outperforms \texttt{CatBoost\_ref}
  in Phase~II (41.70 vs.\ 41.72\,min MAE; 51.86 vs.\ 52.24\,min RMSE).

  \item \textbf{In-race checkpoints ($C_1$ and $C_2$).} Once observed
  split times become available, tree ensembles consistently surpass LR
  in MAE and RMSE. At $C_1$, \texttt{CatBoost\_ref} improves upon LR
  by 2.34\,min in Phase~I (28.43 vs.\ 30.77\,min) and by 0.98\,min in
  Phase~II (33.45 vs.\ 34.43\,min). At $C_2$, \texttt{CatBoost\_ref}
  achieves the lowest MAE, RMSE, and $P_{90}$ in both phases
  (14.79\,min in Phase~I; 15.79\,min in Phase~II), closely followed by
  RF (15.19 and 15.81\,min) and LR (14.89 and 16.55\,min).
\end{itemize}

\subsection{Cross-Period Patterns and Pre-QC Sensitivity}

Test errors are systematically higher in Phase~II than in Phase~I
across all baselines and models (e.g., $B_0$ MAE increases from 40.03
to 45.35\,min; \texttt{CatBoost\_ref} $C_0$ MAE increases from 36.37
to 41.72\,min), accompanied by a marked shift in prediction bias.
Whereas Phase~I biases are small and mostly positive ($+0.13$ to
$+2.11$\,min for tree models; $+2.68$ to $+7.25$\,min for LR),
Phase~II exhibits negative biases across all models, most severely in
demographic baselines ($B_0$: $-11.97$\,min; $B_1/B_2$ at $C_0$:
$-10.26$\,min) and tree ensembles (\texttt{CatBoost\_ref}: $-5.12$\,min
at $C_0$, $-4.43$\,min at $C_1$), while LR remains nearly unbiased
($-0.47$ to $-0.81$\,min).

Evaluating sensitivity to the 113 Phase~II QC exclusions
($Y_i \le W_i + B_i$) confirms why logical-consistency filtering is
vital for quadratic metrics: prior to excluding those 113 records,
Phase~II $C_2$ LR RMSE was inflated to 37.01\,min (vs.\ 22.53\,min
post-QC) and $B_2$ RMSE was 34.23\,min (vs.\ 30.03\,min post-QC),
while \texttt{CatBoost\_ref} MAE changed marginally (15.82 $\to$
15.79\,min) and RMSE improved slightly (21.13 $\to$ 21.02\,min).

% ---------------------------------------------------------------
% 5. Discussion
% ---------------------------------------------------------------
\section{Discussion}

\subsection{Practical Value of Post-Swim and Post-Bike Updates}

A central operational question motivating this work was whether
updating finish-time forecasts immediately after the swim split ($C_1$)
provides sufficient error reduction before the much longer cycling
stage finishes. Our paired results show that knowing \texttt{SwimTime}
alone cuts MAE by 7.94--8.42\,min for tree-based ensembles
(\texttt{CatBoost\_ref} and RF), with \texttt{CatBoost\_ref} beating
the demographic split-addition baseline $B_2$ by over 6.10\,min (and LR
beating $B_2$ by 3.78--5.13\,min). Because $B_2$ adds $W_i$ to a
demographic median of remaining time, this margin demonstrates that
learned models capture richer predictive relationships among swim
split, athlete/event covariates, and overall finish time than simple
demographic remainder medians. Following the bike split ($C_2$),
MAE falls to roughly 15--16\,min ($P_{90} \approx 31$--33\,min),
providing practical real-time guidance for pacing targets and course
safety operations during the concluding half-marathon.

\subsection{Bias Dynamics Under Cross-Source and Temporal Shifts}

A striking empirical pattern in Phase~II is the divergence in bias
between linear and tree-based models. Negative biases in $B_0$
($-11.97$\,min) and $B_1$ ($-10.26$\,min) indicate that 2021--2023
athletes in the Endurance-Data partition finished roughly 10--12\,min
slower on average than the 2004--2018 training medians.
Because tree-based models (RF and \texttt{CatBoost\_ref}) cannot
extrapolate monotonic year trends beyond the 2018 training boundary,
their $C_0$ and $C_1$ predictions inherit negative biases ($-4.43$ to
$-5.12$\,min). In contrast, OLS Linear Regression allows unbounded
linear extrapolation of continuous covariates such as \texttt{EventYear};
combined with the positive $+5.48$ to $+7.25$\,min bias already
exhibited by LR in Phase~I, this linear behavior is consistent with
why LR yields near-zero mean bias ($-0.47$\,min at $C_0$; $-0.50$\,min
at $C_1$) and matches \texttt{CatBoost\_ref} MAE at $C_0$ (41.70 vs.\
41.72\,min) in Phase~II. Consistent with our temporal protocol design, this shift
cannot be attributed solely to biological performance drift, as
changes in source coverage and post-2020 event composition between
Kaggle and Endurance-Data may co-occur.

\subsection{Limitations and Planned Work}

Several limitations of the present empirical study define clear next
steps:
\begin{enumerate}[leftmargin=*]
  \item \textbf{Un-tuned reference configuration.}
  \texttt{CatBoost\_ref} was intentionally evaluated with a fixed
  120-iteration, single-seed (seed~42) baseline; multi-seed stability
  runs and validation-fold hyperparameter tuning remain to be executed.
  \item \textbf{Venue naming and longitudinal identity.} Raw
  \texttt{eventlocation} strings were used after whitespace trimming
  without a verified cross-year alias dictionary, and individual
  athletes lack persistent cross-year tracking IDs across the merged
  sources.
  \item \textbf{Unfinished RQ2 and RQ3 evaluations.} Formal geographic
  ablation ($F_0$--$F_3$ stratified by seen vs.\ unseen venues),
  Location-based Hierarchical Regression (REML), conformal prediction
  intervals (CQR) on the reserved 2017--2018 / 2019--2020 splits, and
  race-cluster bootstrap confidence intervals for $\Delta\mathrm{MAE}$
  require the complete per-record prediction archive and are reserved
  for the next phase of the research program.
\end{enumerate}

% ---------------------------------------------------------------
% 6. Conclusion
% ---------------------------------------------------------------
\section{Conclusion}

We presented a leakage-free, three-checkpoint supervised regression
study on 1{,}152{,}545 quality-controlled IRONMAN~70.3 finishes across
two out-of-period temporal evaluations. Sequential split updates
monotonically reduce prediction error: incorporating swim time at
$C_1$ reduces \texttt{CatBoost\_ref} MAE by 7.94\,min (Phase~I) and
8.27\,min (Phase~II), while adding bike time at $C_2$ yields a further
13.64\,min and 17.67\,min reduction, reaching 14.79\,min and
15.79\,min MAE and shrinking 90th-percentile error by over 57\%. Across
all in-race checkpoints ($C_1$ and $C_2$), tree-based ensembles
(\texttt{CatBoost\_ref} and RF) outperform the checkpoint-aware
demographic remainder baseline $B_2$ by 5.89--7.36\,min (and LR
outperforms $B_2$ by 3.78--6.60\,min).
These findings establish a verified benchmark for sequential triathlon
forecasting and motivate ongoing extensions into hierarchical location
modeling and conformal uncertainty quantification.

% ---------------------------------------------------------------
% References
% ---------------------------------------------------------------
\bibliographystyle{IEEEtran}
\begin{thebibliography}{10}

\bibitem{thuany2025modeling}
M.~Thuany, D.~Valero, E.~Villiger, P.~Forte, K.~Weiss, M.~S.~Andrade,
P.~T.~Nikolaidis, I.~Cuk, T.~Rosemann, and B.~Knechtle,
``Modeling performance in {IRONMAN}~70.3 age group triathletes,''
\emph{Sports Medicine~-- Open}, vol.~11, no.~1, Art.\ no.~148, 2025.
\newblock \doi{10.1186/s40798-025-00948-0}

\bibitem{nikolaidis2023predicting}
P.~T.~Nikolaidis, D.~Valero, K.~Weiss, E.~Villiger, M.~Thuany,
C.~V.~Sousa, M.~Andrade, and B.~Knechtle,
``Predicting overall performance in {Ironman}~70.3 age group
triathletes through split disciplines,''
\emph{Scientific Reports}, vol.~13, Art.\ no.~11492, 2023.
\newblock \doi{10.1038/s41598-023-38181-y}

\bibitem{romano2019conformalized}
Y.~Romano, E.~Patterson, and E.~J.~Cand\`{e}s,
``Conformalized quantile regression,''
in \emph{Advances in Neural Information Processing Systems (NeurIPS)},
vol.~32, 2019, pp.~3538--3548.

\bibitem{prokhorenkova2018catboost}
L.~Prokhorenkova, G.~Gusev, A.~Vorobev, A.~V.~Dorogush,
and A.~Gulin,
``{CatBoost}: Unbiased boosting with categorical features,''
in \emph{Advances in Neural Information Processing Systems (NeurIPS)},
vol.~31, 2018, pp.~6639--6649.

\bibitem{kaggle_dataset}
David,
``{IRONMAN} 70.3 races 2004--2020,''
Kaggle Dataset; Endurance-Data, ``Endurance-data race results
2021--2023.'' [Online]. Available: \url{https://www.endurance-data.com}

\bibitem{repo_artifact}
T.~M.~Khue, N.~S.~Thu, P.~P.~Thao, and D.~T.~G.~Hoi,
``{IRONMAN}~70.3 finish-time prediction: Code, trained models and
reproducibility artifacts,''
GitHub repository, FPT University Da Nang, 2026. [Online].
Available: \url{https://github.com/kaity121/ironman-70-3-finish-time}

\end{thebibliography}

\end{document}
